%%%%%%%%%%%%%%%%%%%%%%%%%%%%%%%%%%%%%%%%%%%%%%%%%%%%%%%%%%%%
% 1.13 HOMOLOGICAL ALGEBRA
% The Composition of Advanced Mathematics
%%%%%%%%%%%%%%%%%%%%%%%%%%%%%%%%%%%%%%%%%%%%%%%%%%%%%%%%%%%%

\section{Homological Algebra}
\label{sec:homological-algebra}

\prerequisite{
Linear algebra from \cref{sec:linear-algebra}, abstract algebra from
\cref{sec:abstract-algebra}, group theory from \cref{sec:group-theory},
ring theory from \cref{sec:ring-theory}, and module theory from
\cref{sec:module-theory}. Familiarity with category-theoretic language is
helpful but not required for the introductory development.
}

\learningobjectives{
By the end of this section, the reader should be able to:
\begin{itemize}
    \item explain the purpose of homological algebra;
    \item interpret kernels, images, and quotients as measurements of exactness;
    \item define chain complexes and cochain complexes;
    \item explain the condition \(d^2=0\);
    \item compute cycles, boundaries, and homology groups;
    \item recognize exact and short exact sequences;
    \item determine whether elementary sequences are exact;
    \item define chain maps and induced maps on homology;
    \item understand chain homotopy at an introductory level;
    \item define projective, injective, and free resolutions conceptually;
    \item explain the role of derived functors;
    \item recognize the functors \(\operatorname{Ext}\) and \(\operatorname{Tor}\);
    \item interpret long exact sequences;
    \item explain how homological algebra connects algebra with topology and geometry.
\end{itemize}
}

%%%%%%%%%%%%%%%%%%%%%%%%%%%%%%%%%%%%%%%%%%%%%%%%%%%%%%%%%%%%
\subsection{Why Homological Algebra?}
\label{subsec:why-homological-algebra}
%%%%%%%%%%%%%%%%%%%%%%%%%%%%%%%%%%%%%%%%%%%%%%%%%%%%%%%%%%%%

Many algebraic problems involve a sequence of maps rather than a single
object.

A typical pattern is

\[ A\xrightarrow{f}B\xrightarrow{g}C. \]

The relationship between the two maps is often more important than either
map individually.

In particular, one asks whether

\[ \Image(f)=\Ker(g). \]

If these two subobjects agree, then everything sent to zero by \(g\) already
came from \(A\).

Homological algebra develops systematic methods for measuring the failure of
such sequences to be exact.

\begin{conceptbox}
Homological algebra studies algebraic objects by organizing them into
sequences of maps and measuring where exactness fails.

The resulting measurements are called \emph{homology}.
\end{conceptbox}

%%%%%%%%%%%%%%%%%%%%%%%%%%%%%%%%%%%%%%%%%%%%%%%%%%%%%%%%%%%%
\subsection{The Fundamental Pattern}
\label{subsec:fundamental-homological-pattern}
%%%%%%%%%%%%%%%%%%%%%%%%%%%%%%%%%%%%%%%%%%%%%%%%%%%%%%%%%%%%

The central picture is

\begin{center}
\begin{tikzcd}[column sep=large]
A \arrow[r,"f"] & B \arrow[r,"g"] & C
\end{tikzcd}
\end{center}

with the compatibility condition

\[ g\circ f=0. \]

This implies

\[ \Image(f)\subseteq\Ker(g). \]

If equality holds, the sequence is exact at \(B\).

If equality does not hold, the quotient

\[ \frac{\Ker(g)}{\Image(f)} \]

measures the discrepancy.

This quotient is the basic prototype of a homology object.

%%%%%%%%%%%%%%%%%%%%%%%%%%%%%%%%%%%%%%%%%%%%%%%%%%%%%%%%%%%%
\subsection{Sequences of Homomorphisms}
\label{subsec:sequences-homomorphisms}
%%%%%%%%%%%%%%%%%%%%%%%%%%%%%%%%%%%%%%%%%%%%%%%%%%%%%%%%%%%%

Consider a sequence

\[ \cdots\longrightarrow A_{n+1}
\xrightarrow{d_{n+1}}A_n
\xrightarrow{d_n}A_{n-1}
\longrightarrow\cdots. \]

The maps

\[ d_n \]

are often called \emph{differentials} or \emph{boundary maps}.

The letter \(d\) is traditionally associated with differentiation and
boundary operations.

The subscript indicates the degree from which the map begins:

\[ d_n:A_n\to A_{n-1}. \]

%%%%%%%%%%%%%%%%%%%%%%%%%%%%%%%%%%%%%%%%%%%%%%%%%%%%%%%%%%%%
\subsection{Chain Complexes}
\label{subsec:chain-complexes}
%%%%%%%%%%%%%%%%%%%%%%%%%%%%%%%%%%%%%%%%%%%%%%%%%%%%%%%%%%%%

\begin{definitionbox}{Chain Complex}
A \emph{chain complex} is a sequence of modules or abelian groups
\[ \cdots\xrightarrow{d_{n+2}}C_{n+1}
\xrightarrow{d_{n+1}}C_n
\xrightarrow{d_n}C_{n-1}
\xrightarrow{d_{n-1}}\cdots \]
such that
\[ d_n\circ d_{n+1}=0 \]
for every \(n\).
\end{definitionbox}

The condition is often abbreviated

\[ d^2=0. \]

This is read ``\(d\) squared equals zero.''

It does not mean that the maps themselves are zero. It means that applying
two consecutive differentials always produces the zero map.

%%%%%%%%%%%%%%%%%%%%%%%%%%%%%%%%%%%%%%%%%%%%%%%%%%%%%%%%%%%%
\subsection{Visual Structure of a Chain Complex}
\label{subsec:chain-complex-visual}
%%%%%%%%%%%%%%%%%%%%%%%%%%%%%%%%%%%%%%%%%%%%%%%%%%%%%%%%%%%%

\begin{center}
\begin{tikzcd}[column sep=large]
\cdots
\arrow[r] &
C_{n+1}
\arrow[r,"d_{n+1}"] &
C_n
\arrow[r,"d_n"] &
C_{n-1}
\arrow[r] &
\cdots
\end{tikzcd}
\end{center}

The defining condition says that every two-step path is zero:

\[ C_{n+1}\xrightarrow{d_{n+1}}C_n\xrightarrow{d_n}C_{n-1} \]

satisfies

\[ d_nd_{n+1}=0. \]

Therefore,

\[ \Image(d_{n+1})\subseteq\Ker(d_n). \]

This inclusion makes the homology quotient possible.

%%%%%%%%%%%%%%%%%%%%%%%%%%%%%%%%%%%%%%%%%%%%%%%%%%%%%%%%%%%%
\subsection{Cycles}
\label{subsec:cycles}
%%%%%%%%%%%%%%%%%%%%%%%%%%%%%%%%%%%%%%%%%%%%%%%%%%%%%%%%%%%%

\begin{definitionbox}{Cycles}
The elements of
\[ Z_n=\Ker(d_n) \]
are called the \emph{\(n\)-cycles}.
\end{definitionbox}

The notation

\[ Z_n \]

is read ``\(Z\) sub \(n\).''

An element \(z\in C_n\) is a cycle when

\[ d_n(z)=0. \]

Thus cycles are elements that have zero boundary.

%%%%%%%%%%%%%%%%%%%%%%%%%%%%%%%%%%%%%%%%%%%%%%%%%%%%%%%%%%%%
\subsection{Boundaries}
\label{subsec:boundaries}
%%%%%%%%%%%%%%%%%%%%%%%%%%%%%%%%%%%%%%%%%%%%%%%%%%%%%%%%%%%%

\begin{definitionbox}{Boundaries}
The elements of
\[ B_n=\Image(d_{n+1}) \]
are called the \emph{\(n\)-boundaries}.
\end{definitionbox}

An element \(b\in C_n\) is a boundary when there exists

\[ c\in C_{n+1} \]

such that

\[ b=d_{n+1}(c). \]

Since \(d^2=0\),

\[ B_n\subseteq Z_n. \]

Every boundary is therefore automatically a cycle.

%%%%%%%%%%%%%%%%%%%%%%%%%%%%%%%%%%%%%%%%%%%%%%%%%%%%%%%%%%%%
\subsection{Homology}
\label{subsec:homology-definition}
%%%%%%%%%%%%%%%%%%%%%%%%%%%%%%%%%%%%%%%%%%%%%%%%%%%%%%%%%%%%

The distinction between cycles and boundaries produces homology.

\begin{definitionbox}{Homology}
The \emph{\(n\)th homology} of a chain complex \(C_\bullet\) is
\[ H_n(C_\bullet)=\frac{Z_n}{B_n}
=\frac{\Ker(d_n)}{\Image(d_{n+1})}. \]
\end{definitionbox}

The notation

\[ C_\bullet \]

is read ``\(C\) bullet'' and represents the entire chain complex.

The notation

\[ H_n(C_\bullet) \]

is read ``the \(n\)th homology of \(C\) bullet.''

%%%%%%%%%%%%%%%%%%%%%%%%%%%%%%%%%%%%%%%%%%%%%%%%%%%%%%%%%%%%
\subsection{What Homology Measures}
\label{subsec:what-homology-measures}
%%%%%%%%%%%%%%%%%%%%%%%%%%%%%%%%%%%%%%%%%%%%%%%%%%%%%%%%%%%%

Because

\[ B_n\subseteq Z_n, \]

the quotient

\[ Z_n/B_n \]

identifies cycles that differ by a boundary.

A nonzero homology class therefore represents a cycle that is not itself a
boundary.

Schematically,

\begin{center}
\begin{tikzpicture}[
    node distance=14mm,
    box/.style={draw,rounded corners,align=center,inner sep=6pt},
    >=stealth
]
\node[box] (Cn) {\(C_n\)};
\node[box,below left=of Cn] (Zn) {cycles\\\(Z_n=\Ker d_n\)};
\node[box,below=of Zn] (Bn) {boundaries\\\(B_n=\Image d_{n+1}\)};
\node[box,right=of Zn] (Hn) {homology\\\(H_n=Z_n/B_n\)};

\draw[->] (Cn)--(Zn);
\draw[->] (Bn)--node[left]{contained in}(Zn);
\draw[->] (Zn)--node[above]{quotient}(Hn);
\end{tikzpicture}
\end{center}

\begin{importantbox}
Homology measures cycles that cannot be explained as boundaries.
\end{importantbox}

%%%%%%%%%%%%%%%%%%%%%%%%%%%%%%%%%%%%%%%%%%%%%%%%%%%%%%%%%%%%
\subsection{A Small Linear-Algebra Example}
\label{subsec:homology-linear-example}
%%%%%%%%%%%%%%%%%%%%%%%%%%%%%%%%%%%%%%%%%%%%%%%%%%%%%%%%%%%%

Consider

\[ \R^2\xrightarrow{d_2}\R^2\xrightarrow{d_1}\R \]

where

\[ d_2(x,y)=(y,0),\qquad d_1(x,y)=y. \]

First check the chain condition:

\[ d_1(d_2(x,y))=d_1(y,0)=0. \]

Thus

\[ d_1d_2=0. \]

At degree \(1\),

\[ Z_1=\Ker(d_1)=\{(x,0):x\in\R\}. \]

Meanwhile,

\[ B_1=\Image(d_2)=\{(y,0):y\in\R\}. \]

Hence

\[ Z_1=B_1. \]

Therefore,

\[ H_1=Z_1/B_1=0. \]

There are no nontrivial homology classes at degree \(1\).

%%%%%%%%%%%%%%%%%%%%%%%%%%%%%%%%%%%%%%%%%%%%%%%%%%%%%%%%%%%%
\subsection{Worked Example: Computing Homology}
\label{subsec:worked-computing-homology}
%%%%%%%%%%%%%%%%%%%%%%%%%%%%%%%%%%%%%%%%%%%%%%%%%%%%%%%%%%%%

\begin{workedexamplebox}{Homology of a Simple Chain Complex}

Consider

\[ 0\longrightarrow\R
\xrightarrow{d_2}\R^2
\xrightarrow{d_1}\R
\longrightarrow0 \]

with

\[ d_2(t)=(t,0),\qquad d_1(x,y)=y. \]

Compute \(H_1\).

\begin{tabularx}{\textwidth}{@{}>{\raggedright\arraybackslash}p{0.42\textwidth}X@{}}
Find the cycles. &
\(\displaystyle Z_1=\Ker(d_1)=\{(x,0):x\in\R\}\)\\
Find the boundaries. &
\(\displaystyle B_1=\Image(d_2)=\{(t,0):t\in\R\}\)\\
Compare the two spaces. & \(Z_1=B_1\)\\
Form the quotient. & \(H_1=Z_1/B_1=0\)
\end{tabularx}

\begin{answerbox}
\[ \boxed{H_1=0} \]
\end{answerbox}
\end{workedexamplebox}

%%%%%%%%%%%%%%%%%%%%%%%%%%%%%%%%%%%%%%%%%%%%%%%%%%%%%%%%%%%%
\subsection{Exact Sequences}
\label{subsec:exact-sequences}
%%%%%%%%%%%%%%%%%%%%%%%%%%%%%%%%%%%%%%%%%%%%%%%%%%%%%%%%%%%%

\begin{definitionbox}{Exact Sequence}
A sequence
\[ \cdots\longrightarrow A_{n+1}
\xrightarrow{f_{n+1}}A_n
\xrightarrow{f_n}A_{n-1}
\longrightarrow\cdots \]
is \emph{exact at \(A_n\)} if
\[ \Image(f_{n+1})=\Ker(f_n). \]
It is \emph{exact} if it is exact at every object where exactness is
required.
\end{definitionbox}

Exactness means that everything killed by one map came from the preceding
map.

%%%%%%%%%%%%%%%%%%%%%%%%%%%%%%%%%%%%%%%%%%%%%%%%%%%%%%%%%%%%
\subsection{Exactness and Homology}
\label{subsec:exactness-homology}
%%%%%%%%%%%%%%%%%%%%%%%%%%%%%%%%%%%%%%%%%%%%%%%%%%%%%%%%%%%%

For a chain complex,

\[ H_n=\frac{\Ker(d_n)}{\Image(d_{n+1})}. \]

Therefore,

\[ H_n=0 \]

precisely when

\[ \Ker(d_n)=\Image(d_{n+1}). \]

Hence:

\begin{theorem}[Exactness and Vanishing Homology]
A chain complex is exact at \(C_n\) if and only if
\[ H_n(C_\bullet)=0. \]
\end{theorem}

Thus homology directly measures failure of exactness.

%%%%%%%%%%%%%%%%%%%%%%%%%%%%%%%%%%%%%%%%%%%%%%%%%%%%%%%%%%%%
\subsection{Short Exact Sequences}
\label{subsec:short-exact-sequences}
%%%%%%%%%%%%%%%%%%%%%%%%%%%%%%%%%%%%%%%%%%%%%%%%%%%%%%%%%%%%

One of the most important configurations in algebra is

\[ 0\longrightarrow A\xrightarrow{f}B\xrightarrow{g}C\longrightarrow0. \]

\begin{definitionbox}{Short Exact Sequence}
A sequence
\[ 0\longrightarrow A\xrightarrow{f}B\xrightarrow{g}C\longrightarrow0 \]
is a \emph{short exact sequence} if it is exact at \(A\), \(B\), and \(C\).
\end{definitionbox}

Exactness implies three important facts.

At \(A\),

\[ \Ker(f)=0, \]

so \(f\) is injective.

At \(B\),

\[ \Image(f)=\Ker(g). \]

At \(C\),

\[ \Image(g)=C, \]

so \(g\) is surjective.

%%%%%%%%%%%%%%%%%%%%%%%%%%%%%%%%%%%%%%%%%%%%%%%%%%%%%%%%%%%%
\subsection{Reading a Short Exact Sequence}
\label{subsec:reading-short-exact-sequence}
%%%%%%%%%%%%%%%%%%%%%%%%%%%%%%%%%%%%%%%%%%%%%%%%%%%%%%%%%%%%

The sequence

\[ 0\longrightarrow A\xrightarrow{f}B\xrightarrow{g}C\longrightarrow0 \]

can be interpreted as saying that \(A\) sits inside \(B\), and the quotient
of \(B\) by that copy of \(A\) is \(C\).

More precisely,

\[ C\cong B/\Image(f). \]

Thus short exact sequences encode extensions of one algebraic object by
another.

\begin{center}
\begin{tikzpicture}[
    node distance=14mm,
    box/.style={draw,rounded corners,align=center,inner sep=5pt},
    >=stealth
]
\node[box] (A) {\(A\)};
\node[box,right=of A] (B) {\(B\)};
\node[box,right=of B] (C) {\(C\)};

\draw[->,thick] (A)--node[above]{injects}(B);
\draw[->,thick] (B)--node[above]{quotient}(C);
\end{tikzpicture}
\end{center}

%%%%%%%%%%%%%%%%%%%%%%%%%%%%%%%%%%%%%%%%%%%%%%%%%%%%%%%%%%%%
\subsection{A Standard Short Exact Sequence}
\label{subsec:standard-short-exact-sequence}
%%%%%%%%%%%%%%%%%%%%%%%%%%%%%%%%%%%%%%%%%%%%%%%%%%%%%%%%%%%%

Let \(n\geq1\). Consider

\[ 0\longrightarrow\Z
\xrightarrow{\times n}\Z
\xrightarrow{\pi}\Z/n\Z
\longrightarrow0. \]

The map

\[ \times n:\Z\to\Z \]

means multiplication by \(n\):

\[ k\mapsto nk. \]

The map

\[ \pi:\Z\to\Z/n\Z \]

is the quotient map

\[ k\mapsto[k]_n. \]

We have

\[ \Ker(\pi)=n\Z=\Image(\times n). \]

Therefore the sequence is exact.

%%%%%%%%%%%%%%%%%%%%%%%%%%%%%%%%%%%%%%%%%%%%%%%%%%%%%%%%%%%%
\subsection{Worked Example: Testing Exactness}
\label{subsec:worked-testing-exactness}
%%%%%%%%%%%%%%%%%%%%%%%%%%%%%%%%%%%%%%%%%%%%%%%%%%%%%%%%%%%%

\begin{workedexamplebox}{A Short Exact Sequence}

Determine whether

\[ 0\longrightarrow\R
\xrightarrow{f}\R^2
\xrightarrow{g}\R
\longrightarrow0 \]

is exact when

\[ f(t)=(t,0),\qquad g(x,y)=y. \]

\begin{tabularx}{\textwidth}{@{}>{\raggedright\arraybackslash}p{0.42\textwidth}X@{}}
Check \(f\) is injective. & \(f(t)=0\Rightarrow t=0\)\\
Find the image of \(f\). &
\(\displaystyle \Image(f)=\{(t,0):t\in\R\}\)\\
Find the kernel of \(g\). &
\(\displaystyle \Ker(g)=\{(x,0):x\in\R\}\)\\
Compare. & \(\Image(f)=\Ker(g)\)\\
Check \(g\) is surjective. & \(g(0,y)=y\)
\end{tabularx}

\begin{answerbox}
\[ \boxed{\text{The sequence is exact.}} \]
\end{answerbox}
\end{workedexamplebox}

%%%%%%%%%%%%%%%%%%%%%%%%%%%%%%%%%%%%%%%%%%%%%%%%%%%%%%%%%%%%
\subsection{Split Exact Sequences}
\label{subsec:split-exact-sequences}
%%%%%%%%%%%%%%%%%%%%%%%%%%%%%%%%%%%%%%%%%%%%%%%%%%%%%%%%%%%%

A short exact sequence

\[ 0\to A\xrightarrow{f}B\xrightarrow{g}C\to0 \]

is especially simple when \(B\) can be decomposed into pieces corresponding
to \(A\) and \(C\).

\begin{definitionbox}{Split Exact Sequence}
A short exact sequence
\[ 0\to A\xrightarrow{f}B\xrightarrow{g}C\to0 \]
is \emph{split} if, under suitable equivalent conditions, there exists a map
\[ s:C\to B \]
such that
\[ g\circ s=\identity_C. \]
\end{definitionbox}

The map \(s\) is called a \emph{section} or \emph{splitting map}.

For modules, a split short exact sequence gives

\[ B\cong A\oplus C. \]

%%%%%%%%%%%%%%%%%%%%%%%%%%%%%%%%%%%%%%%%%%%%%%%%%%%%%%%%%%%%
\subsection{Complexes Versus Exact Sequences}
\label{subsec:complexes-vs-exact}
%%%%%%%%%%%%%%%%%%%%%%%%%%%%%%%%%%%%%%%%%%%%%%%%%%%%%%%%%%%%

Every exact sequence is a chain complex because

\[ \Image(d_{n+1})=\Ker(d_n) \]

implies

\[ \Image(d_{n+1})\subseteq\Ker(d_n), \]

and therefore

\[ d_nd_{n+1}=0. \]

The converse is not true.

A chain complex requires only

\[ \Image(d_{n+1})\subseteq\Ker(d_n). \]

Exactness requires equality.

\begin{center}
\begin{tikzpicture}[
    node distance=12mm,
    box/.style={draw,rounded corners,align=center,inner sep=6pt},
    >=stealth
]
\node[box] (chain) {Chain complex\\\(\Image d_{n+1}\subseteq\Ker d_n\)};
\node[box,below=of chain] (exact) {Exact at \(C_n\)\\\(\Image d_{n+1}=\Ker d_n\)};

\draw[->,thick] (exact)--node[right]{stronger condition}(chain);
\end{tikzpicture}
\end{center}

%%%%%%%%%%%%%%%%%%%%%%%%%%%%%%%%%%%%%%%%%%%%%%%%%%%%%%%%%%%%
\subsection{Cochain Complexes}
\label{subsec:cochain-complexes}
%%%%%%%%%%%%%%%%%%%%%%%%%%%%%%%%%%%%%%%%%%%%%%%%%%%%%%%%%%%%

A closely related convention reverses the indexing direction.

\begin{definitionbox}{Cochain Complex}
A \emph{cochain complex} is a sequence
\[ \cdots\longrightarrow C^{n-1}
\xrightarrow{d^{n-1}}C^n
\xrightarrow{d^n}C^{n+1}
\longrightarrow\cdots \]
such that
\[ d^{n+1}\circ d^n=0. \]
\end{definitionbox}

Superscripts are used rather than subscripts.

The corresponding quotient is called \emph{cohomology}:

\[ H^n(C^\bullet)=\frac{\Ker(d^n)}{\Image(d^{n-1})}. \]

The notation is read ``the \(n\)th cohomology of \(C\) bullet.''

%%%%%%%%%%%%%%%%%%%%%%%%%%%%%%%%%%%%%%%%%%%%%%%%%%%%%%%%%%%%
\subsection{Homology Versus Cohomology}
\label{subsec:homology-vs-cohomology}
%%%%%%%%%%%%%%%%%%%%%%%%%%%%%%%%%%%%%%%%%%%%%%%%%%%%%%%%%%%%

The two conventions can be compared as follows.

\begin{center}
\begin{tabularx}{0.94\textwidth}{@{}lXX@{}}
\toprule
 & \textbf{Homology} & \textbf{Cohomology}\\
\midrule
Index & lower \(C_n\) & upper \(C^n\)\\
Direction & degree decreases & degree increases\\
Maps & \(d_n:C_n\to C_{n-1}\) & \(d^n:C^n\to C^{n+1}\)\\
Invariant & \(H_n\) & \(H^n\)\\
Basic quotient &
\(\Ker d_n/\Image d_{n+1}\) &
\(\Ker d^n/\Image d^{n-1}\)\\
\bottomrule
\end{tabularx}
\end{center}

The distinction becomes especially important in topology, geometry, and
derived functor theory.

%%%%%%%%%%%%%%%%%%%%%%%%%%%%%%%%%%%%%%%%%%%%%%%%%%%%%%%%%%%%
\subsection{Chain Maps}
\label{subsec:chain-maps}
%%%%%%%%%%%%%%%%%%%%%%%%%%%%%%%%%%%%%%%%%%%%%%%%%%%%%%%%%%%%

Homomorphisms between individual modules are not enough when working with
chain complexes.

The maps must respect the differentials.

\begin{definitionbox}{Chain Map}
Let \(C_\bullet\) and \(D_\bullet\) be chain complexes. A \emph{chain map}
\[ f_\bullet:C_\bullet\to D_\bullet \]
is a collection of maps
\[ f_n:C_n\to D_n \]
such that
\[ d_n^D\circ f_n=f_{n-1}\circ d_n^C \]
for every \(n\).
\end{definitionbox}

Equivalently, the following diagram commutes.

\begin{center}
\begin{tikzcd}[column sep=large,row sep=large]
C_n \arrow[r,"d_n^C"] \arrow[d,"f_n"'] &
C_{n-1} \arrow[d,"f_{n-1}"]\\
D_n \arrow[r,"d_n^D"'] &
D_{n-1}
\end{tikzcd}
\end{center}

%%%%%%%%%%%%%%%%%%%%%%%%%%%%%%%%%%%%%%%%%%%%%%%%%%%%%%%%%%%%
\subsection{Chain Maps Preserve Cycles and Boundaries}
\label{subsec:chain-maps-cycles-boundaries}
%%%%%%%%%%%%%%%%%%%%%%%%%%%%%%%%%%%%%%%%%%%%%%%%%%%%%%%%%%%%

Suppose \(z\in C_n\) is a cycle.

Then

\[ d_n^C(z)=0. \]

Since \(f_\bullet\) is a chain map,

\[ d_n^D(f_n(z))=f_{n-1}(d_n^C(z))=0. \]

Thus \(f_n(z)\) is a cycle.

Now suppose

\[ b=d_{n+1}^C(c) \]

is a boundary.

Then

\[ f_n(b)=f_n(d_{n+1}^C(c))
=d_{n+1}^D(f_{n+1}(c)). \]

Therefore \(f_n(b)\) is a boundary.

%%%%%%%%%%%%%%%%%%%%%%%%%%%%%%%%%%%%%%%%%%%%%%%%%%%%%%%%%%%%
\subsection{Induced Maps on Homology}
\label{subsec:induced-homology-map}
%%%%%%%%%%%%%%%%%%%%%%%%%%%%%%%%%%%%%%%%%%%%%%%%%%%%%%%%%%%%

Because chain maps send cycles to cycles and boundaries to boundaries, they
produce maps between homology groups.

\begin{theorem}[Induced Map on Homology]
A chain map
\[ f_\bullet:C_\bullet\to D_\bullet \]
induces a homomorphism
\[ H_n(f):H_n(C_\bullet)\to H_n(D_\bullet) \]
defined by
\[ H_n(f)([z])=[f_n(z)]. \]
\end{theorem}

The square-bracket notation

\[ [z] \]

denotes the homology class represented by the cycle \(z\).

%%%%%%%%%%%%%%%%%%%%%%%%%%%%%%%%%%%%%%%%%%%%%%%%%%%%%%%%%%%%
\subsection{Functoriality of Homology}
\label{subsec:functoriality-homology}
%%%%%%%%%%%%%%%%%%%%%%%%%%%%%%%%%%%%%%%%%%%%%%%%%%%%%%%%%%%%

Homology respects identities and compositions.

If

\[ \identity_{C_\bullet}:C_\bullet\to C_\bullet \]

is the identity chain map, then

\[ H_n(\identity_{C_\bullet})
=\identity_{H_n(C_\bullet)}. \]

If

\[ C_\bullet\xrightarrow{f}D_\bullet\xrightarrow{g}E_\bullet, \]

then

\[ H_n(g\circ f)=H_n(g)\circ H_n(f). \]

Thus homology behaves as a functor.

This categorical viewpoint becomes central in advanced homological algebra.

%%%%%%%%%%%%%%%%%%%%%%%%%%%%%%%%%%%%%%%%%%%%%%%%%%%%%%%%%%%%
\subsection{Chain Homotopy}
\label{subsec:chain-homotopy}
%%%%%%%%%%%%%%%%%%%%%%%%%%%%%%%%%%%%%%%%%%%%%%%%%%%%%%%%%%%%

Two chain maps may be different while still producing the same map on
homology.

The appropriate relation is called \emph{chain homotopy}.

\begin{definitionbox}{Chain Homotopy}
Two chain maps
\[ f,g:C_\bullet\to D_\bullet \]
are \emph{chain homotopic} if there exist maps
\[ h_n:C_n\to D_{n+1} \]
such that
\[ f_n-g_n=d_{n+1}^Dh_n+h_{n-1}d_n^C. \]
\end{definitionbox}

The maps \(h_n\) form the chain homotopy.

%%%%%%%%%%%%%%%%%%%%%%%%%%%%%%%%%%%%%%%%%%%%%%%%%%%%%%%%%%%%
\subsection{Homotopic Maps Induce the Same Homology Map}
\label{subsec:homotopic-same-homology}
%%%%%%%%%%%%%%%%%%%%%%%%%%%%%%%%%%%%%%%%%%%%%%%%%%%%%%%%%%%%

Suppose \(f\) and \(g\) are chain homotopic and \(z\in C_n\) is a cycle.

Since

\[ d_n^C(z)=0, \]

we obtain

\[ f_n(z)-g_n(z)=d_{n+1}^D(h_n(z)). \]

Thus \(f_n(z)\) and \(g_n(z)\) differ by a boundary.

Therefore they represent the same homology class:

\[ [f_n(z)]=[g_n(z)]. \]

Hence

\begin{theorem}[Chain-Homotopy Invariance]
If chain maps \(f\) and \(g\) are chain homotopic, then
\[ H_n(f)=H_n(g) \]
for every \(n\).
\end{theorem}

%%%%%%%%%%%%%%%%%%%%%%%%%%%%%%%%%%%%%%%%%%%%%%%%%%%%%%%%%%%%
\subsection{Free Modules Revisited}
\label{subsec:free-modules-homological}
%%%%%%%%%%%%%%%%%%%%%%%%%%%%%%%%%%%%%%%%%%%%%%%%%%%%%%%%%%%%

A free \(R\)-module has a basis and is isomorphic to a direct sum of copies
of \(R\).

For example,

\[ R^n \]

is free.

Free modules are particularly useful because maps out of them can be
constructed by specifying images of basis elements.

Homological algebra often replaces a complicated module by a sequence of
free modules.

%%%%%%%%%%%%%%%%%%%%%%%%%%%%%%%%%%%%%%%%%%%%%%%%%%%%%%%%%%%%
\subsection{Projective Modules}
\label{subsec:projective-modules-homological}
%%%%%%%%%%%%%%%%%%%%%%%%%%%%%%%%%%%%%%%%%%%%%%%%%%%%%%%%%%%%

Projective modules generalize free modules.

One useful characterization is based on lifting maps.

\begin{definitionbox}{Projective Module}
An \(R\)-module \(P\) is \emph{projective} if whenever
\[ M\xrightarrow{\pi}N\to0 \]
is surjective and
\[ f:P\to N \]
is a homomorphism, there exists
\[ \widetilde f:P\to M \]
such that
\[ \pi\circ\widetilde f=f. \]
\end{definitionbox}

The symbol

\[ \widetilde f \]

is read ``\(f\) tilde.''

It represents a lift of \(f\).

%%%%%%%%%%%%%%%%%%%%%%%%%%%%%%%%%%%%%%%%%%%%%%%%%%%%%%%%%%%%
\subsection{The Projective Lifting Property}
\label{subsec:projective-lifting-visual}
%%%%%%%%%%%%%%%%%%%%%%%%%%%%%%%%%%%%%%%%%%%%%%%%%%%%%%%%%%%%

The lifting condition is represented by

\begin{center}
\begin{tikzcd}[column sep=large,row sep=large]
& P \arrow[d,"f"] \arrow[dl,dashed,"\widetilde f"']\\
M \arrow[r,two heads,"\pi"'] & N \arrow[r] & 0
\end{tikzcd}
\end{center}

The dashed arrow is the desired lift.

Every free module is projective.

Not every projective module is necessarily free.

%%%%%%%%%%%%%%%%%%%%%%%%%%%%%%%%%%%%%%%%%%%%%%%%%%%%%%%%%%%%
\subsection{Injective Modules}
\label{subsec:injective-modules-homological}
%%%%%%%%%%%%%%%%%%%%%%%%%%%%%%%%%%%%%%%%%%%%%%%%%%%%%%%%%%%%

Injective modules satisfy a dual extension property.

\begin{definitionbox}{Injective Module}
An \(R\)-module \(I\) is \emph{injective} if whenever
\[ 0\to A\xrightarrow{\iota}B \]
is injective and
\[ f:A\to I \]
is a homomorphism, there exists
\[ \widetilde f:B\to I \]
such that
\[ \widetilde f\circ\iota=f. \]
\end{definitionbox}

The map originally defined on \(A\) can therefore be extended to \(B\).

%%%%%%%%%%%%%%%%%%%%%%%%%%%%%%%%%%%%%%%%%%%%%%%%%%%%%%%%%%%%
\subsection{Projective Versus Injective}
\label{subsec:projective-vs-injective}
%%%%%%%%%%%%%%%%%%%%%%%%%%%%%%%%%%%%%%%%%%%%%%%%%%%%%%%%%%%%

The two ideas are dual.

\begin{center}
\begin{tabularx}{0.94\textwidth}{@{}lXX@{}}
\toprule
 & \textbf{Projective} & \textbf{Injective}\\
\midrule
Basic behavior & lift through surjections & extend across injections\\
Typical direction & maps out of \(P\) & maps into \(I\)\\
Common use & projective resolutions & injective resolutions\\
\bottomrule
\end{tabularx}
\end{center}

%%%%%%%%%%%%%%%%%%%%%%%%%%%%%%%%%%%%%%%%%%%%%%%%%%%%%%%%%%%%
\subsection{Resolutions}
\label{subsec:resolutions}
%%%%%%%%%%%%%%%%%%%%%%%%%%%%%%%%%%%%%%%%%%%%%%%%%%%%%%%%%%%%

A resolution replaces an algebraic object by an exact sequence of simpler
objects.

\begin{definitionbox}{Projective Resolution}
A \emph{projective resolution} of an \(R\)-module \(M\) is an exact sequence
\[ \cdots\longrightarrow P_2\longrightarrow P_1
\longrightarrow P_0\longrightarrow M\longrightarrow0 \]
where every \(P_n\) is projective.
\end{definitionbox}

If each \(P_n\) is free, the sequence is called a \emph{free resolution}.

%%%%%%%%%%%%%%%%%%%%%%%%%%%%%%%%%%%%%%%%%%%%%%%%%%%%%%%%%%%%
\subsection{Visualizing a Resolution}
\label{subsec:resolution-visual}
%%%%%%%%%%%%%%%%%%%%%%%%%%%%%%%%%%%%%%%%%%%%%%%%%%%%%%%%%%%%

\begin{center}
\begin{tikzpicture}[
    node distance=8mm and 10mm,
    box/.style={draw,rounded corners,inner sep=5pt},
    >=stealth
]
\node[box] (P2) {\(P_2\)};
\node[box,right=of P2] (P1) {\(P_1\)};
\node[box,right=of P1] (P0) {\(P_0\)};
\node[box,right=of P0] (M) {\(M\)};
\node[right=of M] (zero) {\(0\)};

\draw[->] (P2)--(P1);
\draw[->] (P1)--(P0);
\draw[->] (P0)--(M);
\draw[->] (M)--(zero);

\node[below=5mm of P1] {projective or free modules};
\end{tikzpicture}
\end{center}

The module \(M\) may be complicated, but the \(P_n\) are chosen to have
better algebraic properties.

%%%%%%%%%%%%%%%%%%%%%%%%%%%%%%%%%%%%%%%%%%%%%%%%%%%%%%%%%%%%
\subsection{A Free Resolution of \(\Z/n\Z\)}
\label{subsec:resolution-zn}
%%%%%%%%%%%%%%%%%%%%%%%%%%%%%%%%%%%%%%%%%%%%%%%%%%%%%%%%%%%%

As a \(\Z\)-module, \(\Z/n\Z\) has the free resolution

\[ 0\longrightarrow\Z
\xrightarrow{\times n}\Z
\longrightarrow\Z/n\Z
\longrightarrow0. \]

Both copies of \(\Z\) are free \(\Z\)-modules.

The map

\[ \Z\to\Z/n\Z \]

is the natural quotient map.

This small resolution is extremely useful when computing \(\operatorname{Tor}\)
and \(\operatorname{Ext}\).

%%%%%%%%%%%%%%%%%%%%%%%%%%%%%%%%%%%%%%%%%%%%%%%%%%%%%%%%%%%%
\subsection{Why Resolutions Matter}
\label{subsec:why-resolutions-matter}
%%%%%%%%%%%%%%%%%%%%%%%%%%%%%%%%%%%%%%%%%%%%%%%%%%%%%%%%%%%%

Suppose one wants to apply a functor to a module \(M\).

The functor may fail to preserve exact sequences.

Instead of applying it directly to \(M\), one may:

\begin{enumerate}
    \item construct a projective or injective resolution;
    \item apply the functor to the resolution;
    \item compute the resulting homology or cohomology.
\end{enumerate}

The resulting invariants measure the failure of exactness.

This is the origin of \emph{derived functors}.

%%%%%%%%%%%%%%%%%%%%%%%%%%%%%%%%%%%%%%%%%%%%%%%%%%%%%%%%%%%%
\subsection{Functors and Exactness}
\label{subsec:functors-exactness}
%%%%%%%%%%%%%%%%%%%%%%%%%%%%%%%%%%%%%%%%%%%%%%%%%%%%%%%%%%%%

Suppose

\[ 0\to A\to B\to C\to0 \]

is exact.

After applying a functor \(F\), one obtains

\[ F(A)\to F(B)\to F(C). \]

The new sequence may no longer be exact everywhere.

A functor that preserves every short exact sequence is called an
\emph{exact functor}.

Many important functors preserve only part of an exact sequence.

%%%%%%%%%%%%%%%%%%%%%%%%%%%%%%%%%%%%%%%%%%%%%%%%%%%%%%%%%%%%
\subsection{Left and Right Exact Functors}
\label{subsec:left-right-exact}
%%%%%%%%%%%%%%%%%%%%%%%%%%%%%%%%%%%%%%%%%%%%%%%%%%%%%%%%%%%%

A covariant functor \(F\) is \emph{left exact} if

\[ 0\to A\to B\to C \]

exact implies

\[ 0\to F(A)\to F(B)\to F(C) \]

is exact.

It is \emph{right exact} if

\[ A\to B\to C\to0 \]

exact implies

\[ F(A)\to F(B)\to F(C)\to0 \]

is exact.

Derived functors measure the part of exactness that is lost.

%%%%%%%%%%%%%%%%%%%%%%%%%%%%%%%%%%%%%%%%%%%%%%%%%%%%%%%%%%%%
\subsection{The Hom Functor}
\label{subsec:hom-functor-homological}
%%%%%%%%%%%%%%%%%%%%%%%%%%%%%%%%%%%%%%%%%%%%%%%%%%%%%%%%%%%%

For a fixed \(R\)-module \(M\), one can consider

\[ \operatorname{Hom}_R(M,-). \]

This is a covariant functor in the second argument.

It is generally left exact but not necessarily right exact.

Its right derived functors produce

\[ \operatorname{Ext}_R^n(M,N). \]

%%%%%%%%%%%%%%%%%%%%%%%%%%%%%%%%%%%%%%%%%%%%%%%%%%%%%%%%%%%%
\subsection{Ext}
\label{subsec:ext}
%%%%%%%%%%%%%%%%%%%%%%%%%%%%%%%%%%%%%%%%%%%%%%%%%%%%%%%%%%%%

The notation

\[ \operatorname{Ext}_R^n(M,N) \]

is read

``Ext over \(R\), degree \(n\), of \(M\) and \(N\).''

The name \emph{Ext} comes from \emph{extension}.

At a conceptual level, Ext measures certain failures of exactness of the
Hom functor.

\begin{conceptbox}
Very roughly,

\[ \boxed{\operatorname{Ext}\text{ measures extension phenomena and the
failure of Hom to be fully exact.}} \]
\end{conceptbox}

In degree zero,

\[ \operatorname{Ext}_R^0(M,N)\cong\operatorname{Hom}_R(M,N). \]

%%%%%%%%%%%%%%%%%%%%%%%%%%%%%%%%%%%%%%%%%%%%%%%%%%%%%%%%%%%%
\subsection{Extensions and \(\operatorname{Ext}^1\)}
\label{subsec:ext1-extensions}
%%%%%%%%%%%%%%%%%%%%%%%%%%%%%%%%%%%%%%%%%%%%%%%%%%%%%%%%%%%%

One reason for the name Ext is that

\[ \operatorname{Ext}_R^1(M,N) \]

is closely related to equivalence classes of module extensions

\[ 0\to N\to E\to M\to0. \]

A zero class corresponds, in the standard interpretation, to a split
extension.

Thus \(\operatorname{Ext}^1\) records how one module may be assembled as an
extension of another.

%%%%%%%%%%%%%%%%%%%%%%%%%%%%%%%%%%%%%%%%%%%%%%%%%%%%%%%%%%%%
\subsection{The Tensor Product Functor}
\label{subsec:tensor-functor-homological}
%%%%%%%%%%%%%%%%%%%%%%%%%%%%%%%%%%%%%%%%%%%%%%%%%%%%%%%%%%%%

For a fixed module \(M\), consider

\[ M\otimes_R -. \]

The tensor symbol

\[ \otimes \]

is read ``tensor.''

The tensor product functor is generally right exact but need not be left
exact.

Its left derived functors produce Tor.

%%%%%%%%%%%%%%%%%%%%%%%%%%%%%%%%%%%%%%%%%%%%%%%%%%%%%%%%%%%%
\subsection{Tor}
\label{subsec:tor}
%%%%%%%%%%%%%%%%%%%%%%%%%%%%%%%%%%%%%%%%%%%%%%%%%%%%%%%%%%%%

The notation

\[ \operatorname{Tor}_n^R(M,N) \]

is read

``Tor over \(R\), degree \(n\), of \(M\) and \(N\).''

At degree zero,

\[ \operatorname{Tor}_0^R(M,N)\cong M\otimes_RN. \]

Higher Tor groups measure failure of the tensor product to preserve
injectivity and exactness on the left.

\begin{conceptbox}
Very roughly,

\[ \boxed{\operatorname{Tor}\text{ measures the failure of tensor products
to be fully exact.}} \]
\end{conceptbox}

%%%%%%%%%%%%%%%%%%%%%%%%%%%%%%%%%%%%%%%%%%%%%%%%%%%%%%%%%%%%
\subsection{Computing Tor from a Resolution}
\label{subsec:computing-tor}
%%%%%%%%%%%%%%%%%%%%%%%%%%%%%%%%%%%%%%%%%%%%%%%%%%%%%%%%%%%%

Suppose

\[ \cdots\to P_1\to P_0\to M\to0 \]

is a projective resolution of \(M\).

Tensor the projective part with \(N\):

\[ \cdots\to P_1\otimes_RN\to P_0\otimes_RN\to0. \]

The homology of this new complex gives

\[ \operatorname{Tor}_n^R(M,N). \]

Schematically,

\begin{center}
\begin{tikzpicture}[
    node distance=12mm,
    box/.style={draw,rounded corners,align=center,inner sep=5pt},
    >=stealth
]
\node[box] (M) {module \(M\)};
\node[box,right=of M] (res) {projective\\resolution};
\node[box,right=of res] (tensor) {tensor with \(N\)};
\node[box,right=of tensor] (hom) {take homology};
\node[box,below=of hom] (tor) {\(\operatorname{Tor}_n^R(M,N)\)};

\draw[->] (M)--(res);
\draw[->] (res)--(tensor);
\draw[->] (tensor)--(hom);
\draw[->] (hom)--(tor);
\end{tikzpicture}
\end{center}

%%%%%%%%%%%%%%%%%%%%%%%%%%%%%%%%%%%%%%%%%%%%%%%%%%%%%%%%%%%%
\subsection{Worked Example: A Basic Tor Computation}
\label{subsec:worked-tor}
%%%%%%%%%%%%%%%%%%%%%%%%%%%%%%%%%%%%%%%%%%%%%%%%%%%%%%%%%%%%

Use the free resolution

\[ 0\to\Z\xrightarrow{\times n}\Z\to\Z/n\Z\to0. \]

Tensor the projective part with \(\Z/m\Z\):

\[ 0\to\Z/m\Z\xrightarrow{\times n}\Z/m\Z\to0. \]

The first homology is the kernel of multiplication by \(n\):

\[ \operatorname{Tor}_1^\Z(\Z/n\Z,\Z/m\Z)
\cong\Ker\left(
\Z/m\Z\xrightarrow{\times n}\Z/m\Z
\right). \]

The result is

\[ \boxed{
\operatorname{Tor}_1^\Z(\Z/n\Z,\Z/m\Z)
\cong\Z/\gcd(n,m)\Z.} \]

This formula illustrates how Tor detects shared torsion.

%%%%%%%%%%%%%%%%%%%%%%%%%%%%%%%%%%%%%%%%%%%%%%%%%%%%%%%%%%%%
\subsection{Computing Ext from a Resolution}
\label{subsec:computing-ext}
%%%%%%%%%%%%%%%%%%%%%%%%%%%%%%%%%%%%%%%%%%%%%%%%%%%%%%%%%%%%

Again let

\[ \cdots\to P_1\to P_0\to M\to0 \]

be a projective resolution.

Apply

\[ \operatorname{Hom}_R(-,N). \]

Because Hom is contravariant in its first argument, the arrows reverse:

\[ 0\to\operatorname{Hom}_R(P_0,N)
\to\operatorname{Hom}_R(P_1,N)
\to\operatorname{Hom}_R(P_2,N)\to\cdots. \]

This is a cochain complex.

Its cohomology gives

\[ \operatorname{Ext}_R^n(M,N). \]

%%%%%%%%%%%%%%%%%%%%%%%%%%%%%%%%%%%%%%%%%%%%%%%%%%%%%%%%%%%%
\subsection{Worked Example: A Basic Ext Computation}
\label{subsec:worked-ext}
%%%%%%%%%%%%%%%%%%%%%%%%%%%%%%%%%%%%%%%%%%%%%%%%%%%%%%%%%%%%

Start with

\[ 0\to\Z\xrightarrow{\times n}\Z\to\Z/n\Z\to0. \]

Apply

\[ \operatorname{Hom}_{\Z}(-,\Z). \]

Since

\[ \operatorname{Hom}_{\Z}(\Z,\Z)\cong\Z, \]

we obtain

\[ 0\longrightarrow\Z\xrightarrow{\times n}\Z\longrightarrow0. \]

The first cohomology is the cokernel:

\[ \operatorname{Ext}_{\Z}^1(\Z/n\Z,\Z)
\cong\Z/n\Z. \]

Therefore,

\begin{answerbox}
\[ \boxed{\operatorname{Ext}_{\Z}^1(\Z/n\Z,\Z)\cong\Z/n\Z.} \]
\end{answerbox}

%%%%%%%%%%%%%%%%%%%%%%%%%%%%%%%%%%%%%%%%%%%%%%%%%%%%%%%%%%%%
\subsection{Long Exact Sequences}
\label{subsec:long-exact-sequences}
%%%%%%%%%%%%%%%%%%%%%%%%%%%%%%%%%%%%%%%%%%%%%%%%%%%%%%%%%%%%

One of the central strengths of homological algebra is that a short exact
sequence of complexes can generate a long exact sequence of homology groups.

Suppose

\[ 0\to A_\bullet\to B_\bullet\to C_\bullet\to0 \]

is a short exact sequence of chain complexes.

Then there is a long exact sequence

\[ \cdots\to H_n(A)
\to H_n(B)
\to H_n(C)
\xrightarrow{\partial_n}
H_{n-1}(A)
\to H_{n-1}(B)
\to\cdots. \]

The symbol

\[ \partial_n \]

uses the partial-boundary symbol \(\partial\), pronounced ``partial,'' and is
read here as ``the connecting homomorphism in degree \(n\).''

%%%%%%%%%%%%%%%%%%%%%%%%%%%%%%%%%%%%%%%%%%%%%%%%%%%%%%%%%%%%
\subsection{The Connecting Homomorphism}
\label{subsec:connecting-homomorphism}
%%%%%%%%%%%%%%%%%%%%%%%%%%%%%%%%%%%%%%%%%%%%%%%%%%%%%%%%%%%%

The map

\[ \partial_n:H_n(C)\to H_{n-1}(A) \]

is called the \emph{connecting homomorphism}.

It links homology groups in adjacent degrees and is responsible for turning
a short exact sequence into a long exact sequence.

Schematically,

\begin{center}
\begin{tikzcd}[column sep=small]
\cdots \arrow[r] &
H_n(A) \arrow[r] &
H_n(B) \arrow[r] &
H_n(C)
\arrow[out=-35,in=145,looseness=1.4,rr,"\partial_n"] &&
H_{n-1}(A) \arrow[r] &
H_{n-1}(B) \arrow[r] &
H_{n-1}(C) \arrow[r] &
\cdots
\end{tikzcd}
\end{center}

%%%%%%%%%%%%%%%%%%%%%%%%%%%%%%%%%%%%%%%%%%%%%%%%%%%%%%%%%%%%
\subsection{Why Long Exact Sequences Matter}
\label{subsec:why-long-exact}
%%%%%%%%%%%%%%%%%%%%%%%%%%%%%%%%%%%%%%%%%%%%%%%%%%%%%%%%%%%%

Suppose two of the three homology theories in a short exact sequence are
known.

The long exact sequence may allow the third to be determined.

Thus exact sequences turn difficult computations into relationships among
simpler computations.

This principle appears repeatedly throughout modern mathematics.

%%%%%%%%%%%%%%%%%%%%%%%%%%%%%%%%%%%%%%%%%%%%%%%%%%%%%%%%%%%%
\subsection{The Snake Lemma}
\label{subsec:snake-lemma}
%%%%%%%%%%%%%%%%%%%%%%%%%%%%%%%%%%%%%%%%%%%%%%%%%%%%%%%%%%%%

A foundational result behind connecting homomorphisms is the Snake Lemma.

Consider a commutative diagram with exact rows:

\begin{center}
\begin{tikzcd}[column sep=large,row sep=large]
0 \arrow[r] &
A' \arrow[r] \arrow[d,"f'"] &
B' \arrow[r] \arrow[d,"g'"] &
C' \arrow[r] \arrow[d,"h'"] &
0\\
0 \arrow[r] &
A \arrow[r] &
B \arrow[r] &
C \arrow[r] &
0
\end{tikzcd}
\end{center}

Under the standard hypotheses, the Snake Lemma produces an exact sequence

\[ 0\to\Ker(f')\to\Ker(g')\to\Ker(h')
\xrightarrow{\delta}
\operatorname{coker}(f')
\to\operatorname{coker}(g')
\to\operatorname{coker}(h')\to0. \]

The Greek letter

\[ \delta \]

is delta, pronounced ``DEL-tuh.''

The map \(\delta\) is a connecting homomorphism.

%%%%%%%%%%%%%%%%%%%%%%%%%%%%%%%%%%%%%%%%%%%%%%%%%%%%%%%%%%%%
\subsection{Cokernels}
\label{subsec:cokernels-homological}
%%%%%%%%%%%%%%%%%%%%%%%%%%%%%%%%%%%%%%%%%%%%%%%%%%%%%%%%%%%%

The Snake Lemma naturally involves cokernels.

\begin{definitionbox}{Cokernel}
For a homomorphism
\[ f:A\to B, \]
the \emph{cokernel} is
\[ \operatorname{coker}(f)=B/\Image(f). \]
\end{definitionbox}

The cokernel measures failure of \(f\) to be surjective.

Indeed,

\[ \operatorname{coker}(f)=0 \]

if and only if \(f\) is surjective.

Compare this with

\[ \Ker(f)=0, \]

which detects injectivity.

%%%%%%%%%%%%%%%%%%%%%%%%%%%%%%%%%%%%%%%%%%%%%%%%%%%%%%%%%%%%
\subsection{Kernel and Cokernel}
\label{subsec:kernel-cokernel-comparison}
%%%%%%%%%%%%%%%%%%%%%%%%%%%%%%%%%%%%%%%%%%%%%%%%%%%%%%%%%%%%

\begin{center}
\begin{tabularx}{0.9\textwidth}{@{}lXX@{}}
\toprule
 & \textbf{Kernel} & \textbf{Cokernel}\\
\midrule
Definition & \(\Ker(f)\) & \(B/\Image(f)\)\\
Detects & failure of injectivity & failure of surjectivity\\
Vanishing means & \(f\) injective & \(f\) surjective\\
\bottomrule
\end{tabularx}
\end{center}

These two constructions are fundamental throughout homological and
categorical algebra.

%%%%%%%%%%%%%%%%%%%%%%%%%%%%%%%%%%%%%%%%%%%%%%%%%%%%%%%%%%%%
\subsection{The Five Lemma}
\label{subsec:five-lemma}
%%%%%%%%%%%%%%%%%%%%%%%%%%%%%%%%%%%%%%%%%%%%%%%%%%%%%%%%%%%%

Another useful result concerns a commutative diagram with exact rows:

\begin{center}
\begin{tikzcd}[column sep=small,row sep=large]
A_1 \arrow[r] \arrow[d,"f_1"] &
A_2 \arrow[r] \arrow[d,"f_2"] &
A_3 \arrow[r] \arrow[d,"f_3"] &
A_4 \arrow[r] \arrow[d,"f_4"] &
A_5 \arrow[d,"f_5"]\\
B_1 \arrow[r] &
B_2 \arrow[r] &
B_3 \arrow[r] &
B_4 \arrow[r] &
B_5
\end{tikzcd}
\end{center}

Under the standard injectivity and surjectivity hypotheses on the surrounding
vertical maps, the Five Lemma allows one to conclude that

\[ f_3 \]

is an isomorphism.

Its importance lies in transferring information through exact diagrams.

%%%%%%%%%%%%%%%%%%%%%%%%%%%%%%%%%%%%%%%%%%%%%%%%%%%%%%%%%%%%
\subsection{Homological Dimension}
\label{subsec:homological-dimension}
%%%%%%%%%%%%%%%%%%%%%%%%%%%%%%%%%%%%%%%%%%%%%%%%%%%%%%%%%%%%

The length of the shortest projective resolution of a module measures how
far that module is from being projective.

\begin{definitionbox}{Projective Dimension}
The \emph{projective dimension} of an \(R\)-module \(M\), denoted
\[ \operatorname{pd}_R(M), \]
is the minimum length of a projective resolution of \(M\), when such a finite
length exists.
\end{definitionbox}

A projective module has

\[ \operatorname{pd}_R(M)=0. \]

If no finite projective resolution exists, its projective dimension is
infinite.

%%%%%%%%%%%%%%%%%%%%%%%%%%%%%%%%%%%%%%%%%%%%%%%%%%%%%%%%%%%%
\subsection{Global Dimension}
\label{subsec:global-dimension}
%%%%%%%%%%%%%%%%%%%%%%%%%%%%%%%%%%%%%%%%%%%%%%%%%%%%%%%%%%%%

The projective dimensions of all modules over a ring can be combined into a
ring invariant.

\begin{definitionbox}{Global Dimension}
The \emph{global dimension} of a ring \(R\) is
\[ \operatorname{gl.dim}(R)
=\sup\{\operatorname{pd}_R(M):M\text{ is an }R\text{-module}\}. \]
\end{definitionbox}

The notation

\[ \sup \]

means supremum and is read ``supremum.''

Global dimension measures the overall homological complexity of a ring.

%%%%%%%%%%%%%%%%%%%%%%%%%%%%%%%%%%%%%%%%%%%%%%%%%%%%%%%%%%%%
\subsection{Homology as an Algebraic Detector}
\label{subsec:homology-detector}
%%%%%%%%%%%%%%%%%%%%%%%%%%%%%%%%%%%%%%%%%%%%%%%%%%%%%%%%%%%%

Homology can be understood as a detector of hidden structure.

\begin{center}
\begin{tikzpicture}[
    node distance=11mm,
    box/.style={draw,rounded corners,align=center,inner sep=5pt},
    >=stealth
]
\node[box] (objects) {Algebraic or\\geometric objects};
\node[box,below=of objects] (complex) {Build a\\chain complex};
\node[box,below=of complex] (cycles) {Find cycles\\\(\Ker d\)};
\node[box,below=of cycles] (boundaries) {Identify boundaries\\\(\Image d\)};
\node[box,below=of boundaries] (homology) {Form quotient\\\(\Ker d/\Image d\)};
\node[box,below=of homology] (invariant) {Homological\\invariant};

\draw[->] (objects)--(complex);
\draw[->] (complex)--(cycles);
\draw[->] (cycles)--(boundaries);
\draw[->] (boundaries)--(homology);
\draw[->] (homology)--(invariant);
\end{tikzpicture}
\end{center}

This same pattern appears in many different branches of mathematics.

%%%%%%%%%%%%%%%%%%%%%%%%%%%%%%%%%%%%%%%%%%%%%%%%%%%%%%%%%%%%
\subsection{Homological Algebra and Topology}
\label{subsec:homological-algebra-topology}
%%%%%%%%%%%%%%%%%%%%%%%%%%%%%%%%%%%%%%%%%%%%%%%%%%%%%%%%%%%%

Homological algebra grew partly from algebraic topology.

A topological space can be converted into a chain complex

\[ \cdots\to C_2(X)\to C_1(X)\to C_0(X)\to0. \]

The resulting homology groups

\[ H_n(X) \]

detect features such as connected components, loops, and higher-dimensional
holes.

The full construction belongs to algebraic topology, but its algebraic
machinery is precisely the machinery developed here.

\begin{applicationbox}
A geometric object can be transformed into algebra:

\[ \text{space}\longrightarrow
\text{chain complex}\longrightarrow
\text{homology groups}. \]

The resulting groups are often easier to compare than the original spaces.
\end{applicationbox}

%%%%%%%%%%%%%%%%%%%%%%%%%%%%%%%%%%%%%%%%%%%%%%%%%%%%%%%%%%%%
\subsection{Homological Algebra and Commutative Algebra}
\label{subsec:homological-commutative-algebra}
%%%%%%%%%%%%%%%%%%%%%%%%%%%%%%%%%%%%%%%%%%%%%%%%%%%%%%%%%%%%

Homological methods are fundamental in commutative algebra.

They are used to study:

\begin{itemize}
    \item projective and free resolutions;
    \item regular sequences;
    \item depth;
    \item projective dimension;
    \item regular local rings;
    \item syzygies;
    \item local cohomology.
\end{itemize}

The canonical development of these ring-theoretic topics belongs to
\cref{sec:commutative-algebra}.

%%%%%%%%%%%%%%%%%%%%%%%%%%%%%%%%%%%%%%%%%%%%%%%%%%%%%%%%%%%%
\subsection{Syzygies}
\label{subsec:syzygies-homological}
%%%%%%%%%%%%%%%%%%%%%%%%%%%%%%%%%%%%%%%%%%%%%%%%%%%%%%%%%%%%

Suppose a module \(M\) has a presentation

\[ R^m\xrightarrow{A}R^n\to M\to0. \]

The relations among the generators of \(M\) are represented by the image of
\(A\).

Relations among those relations form another module.

These higher relations are called \emph{syzygies}.

A free resolution organizes successive layers of syzygies:

\[ \cdots\to F_2\to F_1\to F_0\to M\to0. \]

Thus resolutions systematically record relations, relations among relations,
and so on.

%%%%%%%%%%%%%%%%%%%%%%%%%%%%%%%%%%%%%%%%%%%%%%%%%%%%%%%%%%%%
\subsection{Homological Algebra and Algebraic Geometry}
\label{subsec:homological-algebraic-geometry}
%%%%%%%%%%%%%%%%%%%%%%%%%%%%%%%%%%%%%%%%%%%%%%%%%%%%%%%%%%%%

Modern algebraic geometry uses homological methods extensively.

Examples include:

\begin{itemize}
    \item sheaf cohomology;
    \item derived functors;
    \item derived categories;
    \item Ext groups of sheaves;
    \item spectral sequences;
    \item deformation theory.
\end{itemize}

These ideas allow geometric information to be encoded in algebraic
invariants.

%%%%%%%%%%%%%%%%%%%%%%%%%%%%%%%%%%%%%%%%%%%%%%%%%%%%%%%%%%%%
\subsection{Homological Algebra and Representation Theory}
\label{subsec:homological-representation-theory}
%%%%%%%%%%%%%%%%%%%%%%%%%%%%%%%%%%%%%%%%%%%%%%%%%%%%%%%%%%%%

For representations viewed as modules, Ext groups measure extension
phenomena between representations.

For example,

\[ \operatorname{Ext}^1(M,N) \]

can detect nontrivial ways of combining \(M\) and \(N\) into a larger
representation.

Projective resolutions also allow representations to be studied through
homological dimensions and derived invariants.

%%%%%%%%%%%%%%%%%%%%%%%%%%%%%%%%%%%%%%%%%%%%%%%%%%%%%%%%%%%%
\subsection{Derived Categories: A Preview}
\label{subsec:derived-categories-preview}
%%%%%%%%%%%%%%%%%%%%%%%%%%%%%%%%%%%%%%%%%%%%%%%%%%%%%%%%%%%%

Advanced homological algebra studies chain complexes themselves as primary
objects.

Two complexes that differ in ways invisible to homology are often treated as
equivalent after formally inverting certain maps called
\emph{quasi-isomorphisms}.

\begin{definitionbox}{Quasi-Isomorphism}
A chain map
\[ f:C_\bullet\to D_\bullet \]
is a \emph{quasi-isomorphism} if it induces isomorphisms
\[ H_n(f):H_n(C_\bullet)\xrightarrow{\sim}H_n(D_\bullet) \]
for every \(n\).
\end{definitionbox}

The symbol

\[ \xrightarrow{\sim} \]

is commonly read ``maps isomorphically to.''

Derived categories provide a framework in which quasi-isomorphic complexes
behave as equivalent objects.

%%%%%%%%%%%%%%%%%%%%%%%%%%%%%%%%%%%%%%%%%%%%%%%%%%%%%%%%%%%%
\subsection{Spectral Sequences: A Preview}
\label{subsec:spectral-sequences-preview}
%%%%%%%%%%%%%%%%%%%%%%%%%%%%%%%%%%%%%%%%%%%%%%%%%%%%%%%%%%%%

Some chain complexes are too complicated to compute directly.

A \emph{spectral sequence} organizes a sequence of successive approximations

\[ E_r^{p,q} \]

that, under suitable conditions, converges toward desired homological
information.

The notation

\[ E_r^{p,q} \]

is read ``\(E\) sub \(r\), superscript \(p,q\).''

A typical differential has the form

\[ d_r:E_r^{p,q}\to E_r^{p-r,q+r-1}. \]

Spectral sequences are powerful but technically substantial, so their full
theory belongs to advanced homological algebra.

%%%%%%%%%%%%%%%%%%%%%%%%%%%%%%%%%%%%%%%%%%%%%%%%%%%%%%%%%%%%
\subsection{The Recurring Homological Pattern}
\label{subsec:recurring-homological-pattern}
%%%%%%%%%%%%%%%%%%%%%%%%%%%%%%%%%%%%%%%%%%%%%%%%%%%%%%%%%%%%

Many constructions in homological algebra follow the same general process.

\begin{center}
\begin{tikzpicture}[
    node distance=9mm,
    box/.style={draw,rounded corners,align=center,inner sep=5pt,font=\small},
    >=stealth
]
\node[box] (obj) {Complicated object};
\node[box,below=of obj] (resolve) {Resolve using simpler objects};
\node[box,below=of resolve] (functor) {Apply a functor};
\node[box,below=of functor] (complex) {Obtain a complex};
\node[box,below=of complex] (homology) {Compute homology};
\node[box,below=of homology] (derived) {Derived invariant};

\draw[->] (obj)--(resolve);
\draw[->] (resolve)--(functor);
\draw[->] (functor)--(complex);
\draw[->] (complex)--(homology);
\draw[->] (homology)--(derived);
\end{tikzpicture}
\end{center}

For example:

\[ \text{projective resolution}
\longrightarrow
-\otimes_RN
\longrightarrow
\operatorname{Tor}, \]

while

\[ \text{projective resolution}
\longrightarrow
\operatorname{Hom}_R(-,N)
\longrightarrow
\operatorname{Ext}. \]

%%%%%%%%%%%%%%%%%%%%%%%%%%%%%%%%%%%%%%%%%%%%%%%%%%%%%%%%%%%%
\subsection{Common Errors}
\label{subsec:homological-algebra-common-errors}
%%%%%%%%%%%%%%%%%%%%%%%%%%%%%%%%%%%%%%%%%%%%%%%%%%%%%%%%%%%%

\subsubsection{Assuming \(d^2=0\) Means \(d=0\)}

The condition

\[ d^2=0 \]

means consecutive differentials compose to zero.

Individual differential maps may be highly nontrivial.

%%%%%%%%%%%%%%%%%%%%%%%%%%%%%%%%%%%%%%%%%%%%%%%%%%%%%%%%%%%%
\subsubsection{Confusing Cycles and Boundaries}

Cycles are

\[ Z_n=\Ker(d_n), \]

while boundaries are

\[ B_n=\Image(d_{n+1}). \]

Every boundary is a cycle, but not every cycle must be a boundary.

The difference is exactly what homology measures.

%%%%%%%%%%%%%%%%%%%%%%%%%%%%%%%%%%%%%%%%%%%%%%%%%%%%%%%%%%%%
\subsubsection{Reversing the Homology Quotient}

Homology is

\[ H_n=\frac{\Ker(d_n)}{\Image(d_{n+1})}, \]

not the reverse quotient.

The denominator is contained in the numerator because \(d^2=0\).

%%%%%%%%%%%%%%%%%%%%%%%%%%%%%%%%%%%%%%%%%%%%%%%%%%%%%%%%%%%%
\subsubsection{Confusing a Chain Complex with an Exact Sequence}

A chain complex requires

\[ \Image(d_{n+1})\subseteq\Ker(d_n). \]

Exactness requires

\[ \Image(d_{n+1})=\Ker(d_n). \]

%%%%%%%%%%%%%%%%%%%%%%%%%%%%%%%%%%%%%%%%%%%%%%%%%%%%%%%%%%%%
\subsubsection{Forgetting What Exactness Means at the Ends}

In

\[ 0\to A\xrightarrow{f}B, \]

exactness at \(A\) means \(f\) is injective.

In

\[ B\xrightarrow{g}C\to0, \]

exactness at \(C\) means \(g\) is surjective.

%%%%%%%%%%%%%%%%%%%%%%%%%%%%%%%%%%%%%%%%%%%%%%%%%%%%%%%%%%%%
\subsubsection{Assuming Every Short Exact Sequence Splits}

A short exact sequence need not be split.

Splitting is an additional property.

%%%%%%%%%%%%%%%%%%%%%%%%%%%%%%%%%%%%%%%%%%%%%%%%%%%%%%%%%%%%
\subsubsection{Assuming Every Projective Module Is Free}

Every free module is projective, but projective modules need not always be
free.

The converse depends on the ring and additional hypotheses.

%%%%%%%%%%%%%%%%%%%%%%%%%%%%%%%%%%%%%%%%%%%%%%%%%%%%%%%%%%%%
\subsubsection{Confusing Ext and Tor}

Ext is associated primarily with derived Hom constructions.

Tor is associated primarily with derived tensor-product constructions.

%%%%%%%%%%%%%%%%%%%%%%%%%%%%%%%%%%%%%%%%%%%%%%%%%%%%%%%%%%%%
\subsubsection{Ignoring Variance of Hom}

The functor

\[ \operatorname{Hom}_R(M,N) \]

is contravariant in \(M\) and covariant in \(N\).

Consequently, applying Hom in the first argument reverses arrows.

%%%%%%%%%%%%%%%%%%%%%%%%%%%%%%%%%%%%%%%%%%%%%%%%%%%%%%%%%%%%
\subsubsection{Thinking Homology Is Only Topological}

Homology originated strongly from topology, but homological algebra is an
algebraic framework used throughout algebra, geometry, number theory,
representation theory, and many other fields.

%%%%%%%%%%%%%%%%%%%%%%%%%%%%%%%%%%%%%%%%%%%%%%%%%%%%%%%%%%%%
\subsection{Applications}
\label{subsec:homological-algebra-applications}
%%%%%%%%%%%%%%%%%%%%%%%%%%%%%%%%%%%%%%%%%%%%%%%%%%%%%%%%%%%%

\begin{applicationbox}
\textbf{Algebraic Topology.}
Chain complexes and homology convert geometric information about spaces into
groups and modules.

\medskip

\textbf{Commutative Algebra.}
Free resolutions, projective dimension, Tor, and Ext measure structural
properties of modules and rings.

\medskip

\textbf{Algebraic Geometry.}
Sheaf cohomology and derived functors encode global geometric information
from local algebraic data.

\medskip

\textbf{Representation Theory.}
Ext groups describe extensions of modules and representations, while
projective resolutions reveal deeper representation-theoretic structure.

\medskip

\textbf{Category Theory.}
Homological algebra is naturally formulated through functors, exact
sequences, natural constructions, and derived categories.

\medskip

\textbf{Number Theory.}
Cohomological methods appear throughout modern arithmetic, including group
cohomology and Galois cohomology.

\medskip

\textbf{Differential Geometry.}
Differential forms form cochain complexes whose cohomology captures global
information about manifolds.

\medskip

\textbf{Mathematical Physics.}
Chain complexes and cohomological constructions occur in gauge theory,
topological field theory, and modern formulations of symmetry.
\end{applicationbox}

%%%%%%%%%%%%%%%%%%%%%%%%%%%%%%%%%%%%%%%%%%%%%%%%%%%%%%%%%%%%
\subsection{Exercises}
\label{subsec:homological-algebra-exercises}
%%%%%%%%%%%%%%%%%%%%%%%%%%%%%%%%%%%%%%%%%%%%%%%%%%%%%%%%%%%%

\begin{exercise}[1]
State the definition of a chain complex.
\end{exercise}

\begin{exercise}[1]
Explain the meaning of
\[ d^2=0. \]
\end{exercise}

\begin{exercise}[2]
Prove that
\[ \Image(d_{n+1})\subseteq\Ker(d_n) \]
in every chain complex.
\end{exercise}

\begin{exercise}[1]
Define the cycle module \(Z_n\).
\end{exercise}

\begin{exercise}[1]
Define the boundary module \(B_n\).
\end{exercise}

\begin{exercise}[1]
Define
\[ H_n(C_\bullet). \]
\end{exercise}

\begin{exercise}[2]
Explain why the quotient
\[ Z_n/B_n \]
is well-defined.
\end{exercise}

\begin{exercise}[2]
Let
\[ 0\to\R\xrightarrow{d_2}\R^2\xrightarrow{d_1}\R\to0 \]
where
\[ d_2(t)=(0,t),\qquad d_1(x,y)=x. \]
Compute \(H_1\).
\end{exercise}

\begin{exercise}[2]
Let
\[ d:\R^2\to\R^2 \]
be defined by
\[ d(x,y)=(y,0). \]
Determine whether \(d^2=0\).
\end{exercise}

\begin{exercise}[2]
For the map in the previous exercise, compute \(\Ker(d)\) and
\(\Image(d)\).
\end{exercise}

\begin{exercise}[1]
Define exactness at an object \(B\) in
\[ A\to B\to C. \]
\end{exercise}

\begin{exercise}[2]
Prove that a chain complex is exact at \(C_n\) if and only if \(H_n=0\).
\end{exercise}

\begin{exercise}[1]
Define a short exact sequence.
\end{exercise}

\begin{exercise}[2]
Show that
\[ 0\to A\xrightarrow{f}B \]
is exact at \(A\) if and only if \(f\) is injective.
\end{exercise}

\begin{exercise}[2]
Show that
\[ B\xrightarrow{g}C\to0 \]
is exact at \(C\) if and only if \(g\) is surjective.
\end{exercise}

\begin{exercise}[2]
Verify that
\[ 0\to\Z\xrightarrow{\times n}\Z\to\Z/n\Z\to0 \]
is exact.
\end{exercise}

\begin{exercise}[1]
Define a split short exact sequence.
\end{exercise}

\begin{exercise}[2]
Show that
\[ 0\to A\to A\oplus C\to C\to0 \]
with the natural inclusion and projection is split exact.
\end{exercise}

\begin{exercise}[1]
Define a cochain complex.
\end{exercise}

\begin{exercise}[1]
Define \(H^n(C^\bullet)\).
\end{exercise}

\begin{exercise}[1]
Explain the indexing difference between homology and cohomology.
\end{exercise}

\begin{exercise}[1]
Define a chain map.
\end{exercise}

\begin{exercise}[2]
Prove that a chain map sends cycles to cycles.
\end{exercise}

\begin{exercise}[2]
Prove that a chain map sends boundaries to boundaries.
\end{exercise}

\begin{exercise}[2]
Explain why a chain map induces a well-defined map on homology.
\end{exercise}

\begin{exercise}[2]
Verify that homology respects composition of chain maps.
\end{exercise}

\begin{exercise}[1]
Define chain homotopy.
\end{exercise}

\begin{exercise}[3]
Prove that chain-homotopic maps induce the same map on homology.
\end{exercise}

\begin{exercise}[1]
Define a projective module using the lifting property.
\end{exercise}

\begin{exercise}[2]
Explain why free modules are projective.
\end{exercise}

\begin{exercise}[1]
Define an injective module.
\end{exercise}

\begin{exercise}[1]
Define a projective resolution.
\end{exercise}

\begin{exercise}[1]
Define a free resolution.
\end{exercise}

\begin{exercise}[2]
Explain why
\[ 0\to\Z\xrightarrow{\times n}\Z\to\Z/n\Z\to0 \]
is a free resolution of \(\Z/n\Z\).
\end{exercise}

\begin{exercise}[1]
Explain what it means for a functor to be exact.
\end{exercise}

\begin{exercise}[1]
Explain the difference between left exact and right exact.
\end{exercise}

\begin{exercise}[1]
What derived functors arise from
\[ \operatorname{Hom}_R(M,-)? \]
\end{exercise}

\begin{exercise}[1]
What derived functors arise from
\[ M\otimes_R-? \]
\end{exercise}

\begin{exercise}[1]
State the degree-zero relationships
\[ \operatorname{Ext}_R^0(M,N)\cong\operatorname{Hom}_R(M,N) \]
and
\[ \operatorname{Tor}_0^R(M,N)\cong M\otimes_RN. \]
\end{exercise}

\begin{exercise}[2]
Use the standard free resolution of \(\Z/n\Z\) to compute
\[ \operatorname{Tor}_1^\Z(\Z/n\Z,\Z). \]
\end{exercise}

\begin{exercise}[3]
Use a free resolution to compute
\[ \operatorname{Tor}_1^\Z(\Z/6\Z,\Z/15\Z). \]
\end{exercise}

\begin{exercise}[3]
Compute
\[ \operatorname{Ext}_{\Z}^1(\Z/5\Z,\Z). \]
\end{exercise}

\begin{exercise}[2]
Explain conceptually what
\[ \operatorname{Ext}^1_R(M,N) \]
has to do with short exact sequences.
\end{exercise}

\begin{exercise}[1]
Define the cokernel of a homomorphism.
\end{exercise}

\begin{exercise}[2]
Prove that
\[ \operatorname{coker}(f)=0 \]
if and only if \(f\) is surjective.
\end{exercise}

\begin{exercise}[2]
Explain how kernels and cokernels respectively measure failures of
injectivity and surjectivity.
\end{exercise}

\begin{exercise}[2]
State the form of the long exact sequence in homology arising from
\[ 0\to A_\bullet\to B_\bullet\to C_\bullet\to0. \]
\end{exercise}

\begin{exercise}[1]
What is the role of the connecting homomorphism
\[ \partial_n:H_n(C)\to H_{n-1}(A)? \]
\end{exercise}

\begin{exercise}[2]
State the conclusion of the Snake Lemma in terms of kernels and cokernels.
\end{exercise}

\begin{exercise}[1]
Describe the purpose of the Five Lemma.
\end{exercise}

\begin{exercise}[1]
Define projective dimension.
\end{exercise}

\begin{exercise}[1]
Define global dimension.
\end{exercise}

\begin{exercise}[1]
Define a quasi-isomorphism.
\end{exercise}

\begin{exercise}[2]
Explain why quasi-isomorphisms are important in derived-category theory.
\end{exercise}

\begin{exercise}[2]
Explain the meaning of a syzygy in a free resolution.
\end{exercise}

\begin{exercise}[2]
Describe the general process
\[ \text{resolution}\to\text{functor}\to\text{homology} \]
used to construct derived functors.
\end{exercise}

%%%%%%%%%%%%%%%%%%%%%%%%%%%%%%%%%%%%%%%%%%%%%%%%%%%%%%%%%%%%
\subsection{Section Connections}
\label{subsec:homological-algebra-connections}
%%%%%%%%%%%%%%%%%%%%%%%%%%%%%%%%%%%%%%%%%%%%%%%%%%%%%%%%%%%%

\chapterconnection{
Homological algebra organizes algebraic information into complexes and exact
sequences. The fundamental relation \(d^2=0\) ensures that boundaries are
cycles, allowing the homology quotient
\(\Ker(d_n)/\Image(d_{n+1})\) to measure failure of exactness. Short exact
sequences encode extensions, while chain maps and chain homotopies allow
complexes to be compared structurally. Projective, injective, and free
resolutions replace complicated modules by better-behaved objects. Applying
nonexact functors to these resolutions leads to derived invariants such as
\(\operatorname{Ext}\) and \(\operatorname{Tor}\). Long exact sequences,
connecting homomorphisms, and diagram lemmas provide powerful computational
tools. These ideas form a common algebraic language connecting module
theory, commutative algebra, representation theory, algebraic topology,
algebraic geometry, differential geometry, number theory, and category
theory.
}

%%%%%%%%%%%%%%%%%%%%%%%%%%%%%%%%%%%%%%%%%%%%%%%%%%%%%%%%%%%%
% SECTION SOURCE TRACKING
%%%%%%%%%%%%%%%%%%%%%%%%%%%%%%%%%%%%%%%%%%%%%%%%%%%%%%%%%%%%

%------------------------------------------------------------
% 1.13 Homological Algebra
%------------------------------------------------------------
%
% Sources:
%   - Charles A. Weibel,
%     An Introduction to Homological Algebra.
%     Key: weibel1994homological
%
% Used for:
%   - chain complexes and homology;
%   - exact sequences;
%   - chain maps and chain homotopy;
%   - projective and injective modules;
%   - resolutions;
%   - derived-functor framework;
%   - Ext and Tor;
%   - long exact sequences;
%   - introductory homological dimension;
%   - advanced connections and terminology.
%
% Notes:
%   - Module theory is treated canonically in Section 1.8.
%   - Commutative-algebra applications belong canonically to Section 1.9.
%   - Representation-theoretic applications belong canonically to
%     Section 1.11.
%   - Category-theoretic foundations and categorical formulations of
%     kernels, cokernels, exactness, and derived constructions are developed
%     canonically in Section 1.16.
%   - Topological homology and cohomology receive their canonical treatment
%     in the Algebraic Topology branch.
%   - Derived categories and spectral sequences are previews only and are
%     not intended here as complete advanced treatments.
%   - Visualizations and pedagogical organization are independently
%     developed for this text.
%
\nocite{weibel1994homological}